\documentclass{article}

\usepackage[a4paper,landscape,margin=2cm]{geometry}
\usepackage{../../../../components/mathtex}
\usepackage{colortbl}
\usepackage{array}

\usepackage{fancyhdr}
\pagestyle{fancy}
\fancyhf{}
\fancyhead[L]{DL3 Math-Info}
\fancyhead[C]{Fiche d'entraînement hebdomadaire N°1}
\fancyhead[R]{2026-2027}
\fancyfoot[L]{Ewen Rodrigues de Oliveira}
\fancyfoot[R]{\thepage}

\usepackage{paracol}

\begin{document}
\setlength{\columnsep}{2cm}

\begin{paracol}{2}

\renewcommand{\arraystretch}{1.3}
\noindent
\begin{tabular}{|p{0.6\columnwidth}|>{\centering\arraybackslash}p{0.35\columnwidth}|}
\hline
\rowcolor[RGB]{230,237,247}
\textbf{Matière} & \textbf{Exercices} \\
\hline
Intégration et probabilités & $\emptyset$\\
\hline
Groupes et actions & $\emptyset$\\
\hline
Calcul différentiel & $1,2,3,4,5,6$\\
\hline
Systèmes d'exploitation & $\emptyset$\\
\hline
Algorithmique & $\emptyset$\\
\hline
Programmation fonctionnelle & $\emptyset$\\
\hline
Anglais & $\emptyset$\\
\hline
\end{tabular}

\vspace{1em}

\exercise{1}{Formes hermitiennes et vérifications de base}{
    On considère la fonction $b : \mathbb{C}^2 \times \mathbb{C}^2 \to \mathbb{C}$ définie pour $x=(x_1,x_2)$ et $y=(y_1,y_2)$ par :
    \[
        b(x,y) = 2\overline{x_1}y_1 + i\overline{x_1}y_2 - i\overline{x_2}y_1 + 3\overline{x_2}y_2
    \]
    \begin{enumerate}
        \item Vérifier que $b$ est linéaire à droite (par rapport à $y$) et semi-linéaire à gauche (par rapport à $x$).
        \item Montrer que $b(y,x) = \overline{b(x,y)}$ pour tous $x,y \in \mathbb{C}^2$.
        \item Pour $x=(x_1, x_2) \in \mathbb{C}^2$, exprimer $b(x,x)$ sous la forme d'une somme de carrés de modules.
        \item En déduire que $b$ définit un produit scalaire (hermitien) sur $\mathbb{C}^2$.
    \end{enumerate}
}

\exercise{2}{Identités de polarisation et produit scalaire}{
    Soit $E$ un $\mathbb{C}$-espace vectoriel muni d'un produit scalaire $\langle \cdot, \cdot \rangle$ (linéaire à droite) et de la norme associée $\|\cdot\|$.
    \begin{enumerate}
        \item Développer les expressions $\|x+y\|^2$, $\|x-y\|^2$, $\|x+iy\|^2$ et $\|x-iy\|^2$.
        \item Démontrer l'identité de polarisation complexe :
        \[
            \langle x,y \rangle = \frac{1}{4}\left( \|x+y\|^2 - \|x-y\|^2 - i\|x+i y\|^2 + i\|x-i y\|^2 \right)
        \]
        \item Que devient cette identité si $E$ est un espace préhilbertien réel ?
    \end{enumerate}
}

\switchcolumn

\exercise{3}{Gram-Schmidt dans $\mathbb{C}^3$}{
    Soit $\mathbb{C}^3$ muni de son produit scalaire hermitien canonique $\langle x,y \rangle = \sum_{k=1}^3 \overline{x_k}y_k$.  
    On considère les vecteurs $e_1 = (1, i, 0)$ et $e_2 = (0, 1, i)$.
    \begin{enumerate}
        \item Montrer que la famille $(e_1, e_2)$ est libre dans $\mathbb{C}^3$.
        \item Former une famille orthonormée $(f_1, f_2)$ engendrant le sous-espace $F = \text{Vect}(e_1, e_2)$.
        \item Déterminer un vecteur $f_3 \in \mathbb{C}^3$ non nul tel que $(f_1, f_2, f_3)$ forme une base orthonormée de $\mathbb{C}^3$.
    \end{enumerate}
}


\exercise{4}{Distance à un sous-espace et polynômes}{
    On se place dans $E = \mathbb{R}[X]$ muni du produit scalaire :
    \[
        \langle P, Q \rangle = \int_0^1 P(t)Q(t) \, dt
    \]
    On note $F = \mathbb{R}_1[X] = \text{Vect}(1, X)$.
    \begin{enumerate}
        \item Construire une base orthonormée $(Q_0, Q_1)$ du sous-espace $F$ par l'orthogonalisation de Gram-Schmidt appliquée à la base canonique $(1, X)$.
        \item Calculer la projection orthogonale du polynôme $P(X) = X^2$ sur le sous-espace $F$.
        \item En déduire la valeur de la distance $d(X^2, F) = \inf_{(a,b) \in \mathbb{R}^2} \left( \int_0^1 (t^2 - at - b)^2 \, dt \right)^{1/2}$.
    \end{enumerate}
}

\exercise{5}{Projections dans un espace euclidien}{
    Dans $\mathbb{R}^4$ muni du produit scalaire usuel, on considère le sous-espace vectoriel $F$ d'équations :
    \[
        \begin{cases}
            x + y + z + t = 0 \\
            x - y + z - t = 0
        \end{cases}
    \]
    \begin{enumerate}
        \item Déterminer une base de $F$, puis une base de $F^\perp$. En déduire une base orthonormée de $F$.
        \item Calculer la matrice dans la base canonique de la projection orthogonale $p_F$ sur $F$.
        \item Déterminer la distance du vecteur $v = (1, 2, 3, 4)$ au sous-espace $F$.
    \end{enumerate}
}
\switchcolumn
\newpage

\exercise{6}{L'orthogonal du sous-espace des polynômes}{
    Soit $E = \mathcal{C}^0([0,1], \mathbb{R})$ muni du produit scalaire $\langle f,g \rangle = \int_0^1 f(t)g(t) \, dt$.  
    On note $F = \mathbb{R}[X]$ le sous-espace des fonctions polynomiales restreintes à $[0,1]$.
    \begin{enumerate}
        \item Soit $f \in F^\perp$. Montrer que pour tout $n \in \mathbb{N}$, $\int_0^1 f(t) t^n \, dt = 0$.
        \item En admettant le théorème de Weierstrass (tout élément de $E$ est limite uniforme d'une suite de polynômes sur $[0,1]$), montrer que $\int_0^1 f(t)^2 \, dt = 0$.
        \item En déduire $F^\perp$.
        \item A-t-on $E = F \oplus F^\perp$ ? Que peut-on en conclure sur la dimension infinie ?
    \end{enumerate}
}

\end{paracol}

\end{document}
